\section{FFT}

\subsection{Домашнее задание}
\begin{enumerate}

  \item
    Пусть $p$ --- простое число и $g$ ---
    \href{https://ru.wikipedia.org/wiki/%D0%9F%D0%B5%D1%80%D0%B2%D0%BE%D0%BE%D0%B1%D1%80%D0%B0%D0%B7%D0%BD%D1%8B%D0%B9_%D0%BA%D0%BE%D1%80%D0%B5%D0%BD%D1%8C_(%D1%82%D0%B5%D0%BE%D1%80%D0%B8%D1%8F_%D1%87%D0%B8%D1%81%D0%B5%D0%BB)}{первообразный корень}
    по модулю $p$. Определим \texttt{DFT} размера $p$ обычным образом: $$f_k = \sum_{i=0}^{p-1} x_i \omega_{p}^{k i}$$

    \begin{enumerate}
      \item
        Определим $x^{\prime}_s = x_{g^{-s} \mod p}$, $f^{\prime}_r = f_{g^{r} \mod p}$ и $\Delta_{i} = \omega_{p}^{g^i}$.
        Покажите, что: $$f^{\prime}_{r} = x_{0} + \sum_{s = 0}^{p - 2} x^{\prime}_{s} \Delta_{(r - s) \mod (p - 1)}$$

        \textbf{Решение:}

        \begin{itemize}
          \item По определению:
          $$f'_r = f_{g^r \mod p} = \sum_{i=0}^{p-1} x_i \cdot \omega_p^{(g^r \mod p) \cdot i}$$

          \item Отделим слагаемое $i = 0$:
          $$f'_r = x_0 + \sum_{i=1}^{p-1} x_i \cdot \omega_p^{g^r \cdot i \mod p}$$

          \item Так как $g$ --- первообразный корень по модулю $p$, элементы $\{g^0, g^1, \ldots, g^{p-2}\} \mod p$ --- это в точности $\{1, 2, \ldots, p-1\}$. Поэтому каждое $i \in \{1, \ldots, p-1\}$ единственным образом представимо как $i = g^{-s} \mod p$ для некоторого $s \in \{0, 1, \ldots, p-2\}$.

          \item Подставим $i = g^{-s} \mod p$:
          $$f'_r = x_0 + \sum_{s=0}^{p-2} x_{g^{-s} \mod p} \cdot \omega_p^{g^r \cdot g^{-s} \mod p}$$

          \item Заметим, что $g^r \cdot g^{-s} = g^{r-s}$, и так как порядок $g$ в $(\Z/p\Z)^*$ равен $p - 1$:
          $$g^{r-s} \mod p = g^{(r - s) \mod (p-1)} \mod p$$

          \item Значит:
          $$\omega_p^{g^r \cdot g^{-s} \mod p} = \omega_p^{g^{(r-s) \mod (p-1)}} = \Delta_{(r-s) \mod (p-1)}$$

          \item Подставляя $x'_s = x_{g^{-s} \mod p}$:
          $$f'_r = x_0 + \sum_{s=0}^{p-2} x'_s \cdot \Delta_{(r-s) \mod (p-1)}$$

          ч.т.д.
        \end{itemize}

      \item
        Придумайте, как вычислить $f_k$ за $\O(p \log{p})$ арифметических операций.

        \textbf{Решение:}

        \begin{itemize}
          \item Из пункта (а) следует, что
          $$f'_r = x_0 + \sum_{s=0}^{p-2} x'_s \cdot \Delta_{(r-s) \mod (p-1)}$$

          Сумма $\sum_{s=0}^{p-2} x'_s \cdot \Delta_{(r-s) \mod (p-1)}$ --- это циклическая свёртка двух последовательностей длины $p - 1$: $(x'_0, x'_1, \ldots, x'_{p-2})$ и $(\Delta_0, \Delta_1, \ldots, \Delta_{p-2})$.

          \item Циклическая свёртка длины $N = p - 1$ вычисляется за $\O(N \log N)$:
          \begin{enumerate}
            \item Дополнить обе последовательности нулями до длины $M \geq 2N - 1$, где $M$ --- степень двойки.
            \item Вычислить DFT обеих последовательностей размера $M$.
            \item Поэлементно перемножить результаты.
            \item Вычислить обратное DFT.
            \item Сложить элементы с индексами, совпадающими по модулю $N$ (свернуть обратно до длины $N$).
          \end{enumerate}

          \item После получения свёртки прибавляем $x_0$ к каждому $f'_r$ за $\O(p)$.

          \item Получили все $f'_r = f_{g^r \mod p}$ для $r = 0, \ldots, p-2$, то есть все $f_k$ для $k \in \{1, \ldots, p-1\}$.

          \item Значение $f_0 = \sum_{i=0}^{p-1} x_i$ вычисляем отдельно за $\O(p)$.

          \item Итого: $\O(p \log p)$ арифметических операций.
        \end{itemize}
    \end{enumerate}

  \item
    Придумайте, как свести вычисление \texttt{DFT} последовательности размера $n = n_1 n_2$ к:
    \begin{enumerate}
      \item
        $n_1$ вычислениям \texttt{DFT} от последовательностей размера $n_2$ и
        $\O(n_1^2 n_2)$ дополнительных арифметических операций.
      \item
        $n_1$ вычислениям \texttt{DFT} от последовательностей размера $n_2$,
        $n_2$ вычислениям \texttt{DFT} от последовательностей размера $n_1$ и
        $\O(n_1 n_2)$ дополнительных арифметических операций.
    \end{enumerate}

  \item
    Вам даны две последовательности $a_i$ и $b_i$ длины $n = 2^m$. Вычислите за $\O(n \log{n})$
    арифметических операций последовательность:
    $$c_k = \sum_{i \oplus j = k} a_i b_j$$

    Вам может пригодится отображение $x_i \rightarrow f_i$ определенное следующим образом:
    $$f_k = \sum_{i = 0}^{n - 1} x_i (-1)^{\text{bit\_count}(k \land i)}$$

    \textit{Подсказка: покажите, что:}
    $$(-1)^{\text{bit\_count}(k \land (i \oplus j))} = (-1)^{\text{bit\_count}(k \land i)} \cdot (-1)^{\text{bit\_count}(k \land j)}$$

  \item
    Найдите количество \texttt{AVL}-деревьев высоты $h$ из $n$ вершин по простому модулю $p$ за:
    \begin{enumerate}
      \item $\O(h n \log{n})$.
      \item $\O(h N(h))$. Здесь $N(h)$ --- максимальное количество вершин в \texttt{AVL} дереве высоты $h$.
    \end{enumerate}
    Считайте, что в поле вычетов по модулю $p$ вам известен
    \href{https://ru.wikipedia.org/wiki/%D0%9F%D0%B5%D1%80%D0%B2%D0%BE%D0%BE%D0%B1%D1%80%D0%B0%D0%B7%D0%BD%D1%8B%D0%B9_%D0%BA%D0%BE%D1%80%D0%B5%D0%BD%D1%8C_%D0%B8%D0%B7_%D0%B5%D0%B4%D0%B8%D0%BD%D0%B8%D1%86%D1%8B}{первообразный корень из единицы}
    достаточной степени.

  \item
    Даны строка $s$ длины $n$ и шаблон $p$ длины $m \le n$.
    Алфавит --- \textbf{не} константа.
    \begin{enumerate}
      \item
        Найдите точные вхождения шаблона за $\O(1)$ \texttt{DFT} размера порядка $n$. \\
        \textit{Подсказка: } посчитайте $\sum_{j = 0}^{m - 1} (s_{i+j} - p_j)^2$ для всех $0 \le i \le n - m$.

      \item
        Теперь шаблон может содержать вопросительные знаки. Найдите подстроку $s$,
        точно подходящую под шаблон. $\O(n \log{n})$.

      \item
        Шаблон содержит звёздочки и вопросительные знаки. Проверьте,
        подходит ли строка \textbf{целиком} под шаблон. $\O(n \sqrt{n} \log{n})$.
    \end{enumerate}

\end{enumerate}
